\subsection{Cardinal Arithmetic}

\begin{proposition} \label{prop:cardinal-arithmetic-reference-laws} \hlt{Cardinal Arithmetic Reference Laws}\\
Let $\kappa,\lambda,\mu$ be cardinal numbers.
\begin{enumerate}[label=(\roman*)]
\setlength{\itemsep}{0pt}
\item $\kappa+\lambda=\lambda+\kappa$ and $\kappa\lambda=\lambda\kappa$;
\item $\kappa+(\lambda+\mu)=(\kappa+\lambda)+\mu$ and $\kappa(\lambda\mu)=(\kappa\lambda)\mu$;
\item $\kappa(\lambda+\mu)=\kappa\lambda+\kappa\mu$;
\item $\kappa^{\lambda+\mu}=\kappa^\lambda\kappa^\mu$, $(\kappa^\lambda)^\mu=\kappa^{\lambda\mu}$, and $(\kappa\lambda)^\mu=\kappa^\mu\lambda^\mu$.
\end{enumerate}
\begin{proof}
Choose pairwise disjoint sets $A,B,C$ with $\abs{A}=\kappa$, $\abs{B}=\lambda$, and $\abs{C}=\mu$.
\begin{enumerate}[label=(\roman*)]
\setlength{\itemsep}{0pt}
\item Use $A\cup B\cong B\cup A$ and $A\times B\cong B\times A$.
\item Use $A\cup(B\cup C)\cong(A\cup B)\cup C$ and $A\times(B\times C)\cong(A\times B)\times C$.
\item Use $A\times(B\cup C)\cong(A\times B)\cup(A\times C)$.
\item Use $A^{B\cup C}\cong A^B\times A^C$, $(A^B)^C\cong A^{B\times C}$, and $(A\times B)^C\cong A^C\times B^C$.
\end{enumerate}
Taking cardinalities gives (i)--(iv).\\
\hlt{Structural Commentary}: The laws are canonical bijections.\\
\hlt{Key Insight}: Cardinal arithmetic is set construction up to bijection.
\end{proof}
\end{proposition}

\begin{theorem} \label{thm:infinite-cardinal-absorption-reference} \hlt{Infinite Cardinal Absorption}\\
Let $\kappa,\lambda$ be cardinals and set $\mu\defeq\max\{\kappa,\lambda\}$. If $\mu$ is infinite, then $\kappa+\lambda=\mu$.\\
If also $\kappa,\lambda>0$, then $\kappa\lambda=\mu$. In particular, $\aleph_0+\aleph_0=\aleph_0\cdot\aleph_0=\aleph_0$.
\begin{proof}
For every infinite cardinal $\nu$, the standard square law gives $\nu+\nu=\nu^2=\nu$.\\
Since $\kappa,\lambda\leq\mu$, $\mu\leq \kappa+\lambda\leq\mu+\mu=\mu$, so $\kappa+\lambda=\mu$.\\
If $\kappa,\lambda>0$, choose nonempty representatives $\abs{A}=\kappa$ and $\abs{B}=\lambda$; fixing one point in the smaller factor injects the larger factor into $A\times B$, so $\mu\leq\kappa\lambda$. Also $\kappa\lambda\leq\mu^2=\mu$, so $\kappa\lambda=\mu$.\\
\hlt{Structural Commentary}: Infinite cardinals absorb smaller summands and nonzero factors.\\
\hlt{Key Insight}: Squeeze between $\mu$ and $\mu^2=\mu$.
\end{proof}
\end{theorem}
